\documentclass[11pt]{article}

\usepackage[margin=1in]{geometry}
\usepackage{amsmath, amssymb, bm}
\usepackage{enumitem}
\usepackage{hyperref}

\title{Simple Summary of Operational Space Control Barrier Functions}
\author{Notes for xArm / UF850 Density-Based Planning}
\date{}

\begin{document}
\maketitle

\section{What Problem Is The Paper Solving?}

The paper studies safe real-time control of robot manipulators. The motivating
problem is simple:

\begin{quote}
The end effector may have a desired task, but the whole robot body must remain
safe while executing that task.
\end{quote}

This is directly relevant to our UF850 work. A density controller can generate
a nice end-effector command from point A to point B, but if we simply map that
command through inverse kinematics, the arm links may collide with cylinders,
cuboids, the table, or other obstacles.

The paper's main idea is to combine:

\begin{itemize}[leftmargin=*]
    \item operational-space control, which expresses the task in end-effector
    coordinates;
    \item control barrier functions, which enforce safety constraints;
    \item a task-consistent quadratic program, which changes the command only
    as much as necessary while preserving the task hierarchy.
\end{itemize}

\section{Main Contribution In Simple Terms}

The main contribution is an \textbf{Operational Space Control Barrier Function}
framework, abbreviated OSCBF.

Instead of applying a safety filter only in joint space, the paper designs the
filter so that it understands the task hierarchy:

\[
\text{primary task: end-effector motion}
\]

\[
\text{secondary task: joint posture / null-space behavior}
\]

\[
\text{safety: collision avoidance, joint limits, singularity avoidance, workspace limits}
\]

The controller solves a QP that tries to preserve the operational-space task
while enforcing many safety constraints at once. This is what the authors call
\emph{task consistency}.

\section{Basic CBF Formulation}

Consider a control-affine system

\[
\dot{z} = f(z) + g(z)u,
\]

where:

\begin{itemize}[leftmargin=*]
    \item $z$ is the state,
    \item $u$ is the control input,
    \item $f$ is the natural dynamics,
    \item $g$ describes how the input affects the system.
\end{itemize}

A safety condition is encoded by a scalar function $h(z)$:

\[
\mathcal{C} = \{z : h(z) \ge 0\}.
\]

The set $\mathcal{C}$ is the safe set. If $h(z) < 0$, the system is unsafe.

The CBF condition is:

\[
\dot{h}(z,u) \ge -\alpha(h(z)),
\]

where $\alpha(\cdot)$ is usually a simple increasing function such as

\[
\alpha(h) = \gamma h, \qquad \gamma > 0.
\]

For the control-affine system,

\[
\dot{h}(z,u)
= \frac{\partial h}{\partial z}\dot{z}
= L_f h(z) + L_g h(z)u.
\]

So the CBF constraint becomes

\[
L_f h(z) + L_g h(z)u \ge -\alpha(h(z)).
\]

This is useful because it is usually linear in $u$. Therefore it can be added
as a constraint in a quadratic program.

\section{Safety Filter QP}

A standard CBF safety filter solves:

\[
\begin{aligned}
u^\star
= \arg\min_u \quad & \|u - u_{\mathrm{nom}}\|_2^2 \\
\text{s.t.} \quad
& L_f h(z) + L_g h(z)u \ge -\alpha(h(z)).
\end{aligned}
\]

Here $u_{\mathrm{nom}}$ is the desired, possibly unsafe, control input. The QP
returns $u^\star$, the closest safe input.

In words:

\begin{quote}
Do what the nominal controller wants, unless doing so would violate safety.
\end{quote}

\section{Velocity-Controlled Manipulator Setup}

For our first UF850 implementation, the velocity-controlled version is the most
important.

Let:

\[
z = q,
\qquad
u = \dot{q}.
\]

The reduced-order dynamics are:

\[
\dot{q} = u.
\]

The end-effector velocity is:

\[
\nu = J(q)\dot{q},
\]

where $J(q)$ is the manipulator Jacobian and $\nu$ is the operational-space
twist. For position-only control, we can use the translational Jacobian:

\[
\dot{x} = J_v(q)\dot{q}.
\]

\section{Nominal Operational-Space Command}

The paper computes a nominal end-effector twist using proportional feedback:

\[
\nu_{\mathrm{nom}}
= \nu_{\mathrm{des}}
- K_{p,o}
\begin{bmatrix}
x_p - x_{p,\mathrm{des}} \\
\delta \phi
\end{bmatrix}.
\]

Here:

\begin{itemize}[leftmargin=*]
    \item $x_p$ is the end-effector position,
    \item $x_{p,\mathrm{des}}$ is the desired position,
    \item $\delta \phi$ is orientation error,
    \item $K_{p,o}$ is an operational-space gain.
\end{itemize}

For our density planner, we can replace this nominal command with:

\[
\dot{x}_{\mathrm{nom}}
= u_{\mathrm{density}}(x, x_{\mathrm{goal}}, \text{obstacles}).
\]

That is: density feedback gives the desired end-effector velocity.

\section{Mapping To Joint Velocity}

A naive approach maps the end-effector command to joint velocity using a
Jacobian pseudoinverse:

\[
\dot{q}_{\mathrm{nom}}
= J^\#(q)\nu_{\mathrm{nom}}.
\]

If the robot is redundant, a null-space posture command can be added:

\[
\dot{q}_{\mathrm{nom}}
= J^\#(q)\nu_{\mathrm{nom}}
+ N(q)\dot{q}_0,
\]

where:

\[
N(q) = I - J^\#(q)J(q).
\]

This command is still unsafe by itself. The OSCBF QP filters it.

\section{Task-Consistent OSCBF Objective}

A basic joint-space safety filter would minimize:

\[
\|\dot{q} - \dot{q}_{\mathrm{nom}}\|_2^2.
\]

The paper argues that this can be task-inconsistent. Why? Because the real task
is often the end-effector motion, not matching a particular joint velocity.

Instead, the OSCBF objective penalizes deviations in the task variables:

\[
\|W_o J(q)(\dot{q} - \dot{q}_{\mathrm{nom}})\|_2^2
+
\|W_j N(q)(\dot{q} - \dot{q}_{\mathrm{nom}})\|_2^2.
\]

The first term says:

\[
\text{preserve the operational-space/end-effector command}.
\]

The second term says:

\[
\text{preserve the null-space posture behavior when possible}.
\]

This is the core ``nice'' part of the formulation. The safety filter does not
blindly stay close to joint velocities. It stays close to the behavior that
matters for the task.

\section{Kinematic OSCBF QP}

For a velocity-controlled robot, the OSCBF QP can be written as:

\[
\begin{aligned}
\dot{q}^{\star}
= \arg\min_{\dot{q}} \quad
& \|W_o J(q)(\dot{q} - \dot{q}_{\mathrm{nom}})\|_2^2
 + \|W_j N(q)(\dot{q} - \dot{q}_{\mathrm{nom}})\|_2^2 \\
\text{s.t.} \quad
& \nabla h_k(q)^\top \dot{q}
  \ge -\alpha_k(h_k(q)), \qquad k = 1,\dots,m, \\
& \dot{q}_{\min} \le \dot{q} \le \dot{q}_{\max}.
\end{aligned}
\]

Each $h_k$ is one safety constraint. There can be many of them.

\section{Collision Avoidance Barrier}

The paper uses sphere decompositions for whole-body collision avoidance. The
robot is approximated by spheres attached to links. Obstacles are also
represented as spheres, or decomposed into sphere-like primitives.

For robot sphere $i$ and obstacle sphere $j$:

\[
h_{ij}(q)
= \|p_i(q) - p_{j,\mathrm{obs}}\|_2
- r_i
- r_{j,\mathrm{obs}}.
\]

With a safety margin $d_{\mathrm{safe}}$, use:

\[
h_{ij}(q)
= \|p_i(q) - p_{j,\mathrm{obs}}\|_2
- r_i
- r_{j,\mathrm{obs}}
- d_{\mathrm{safe}}.
\]

Safe means:

\[
h_{ij}(q) \ge 0.
\]

The CBF constraint is:

\[
\nabla h_{ij}(q)^\top \dot{q}
\ge
-\gamma h_{ij}(q).
\]

Intuitively:

\begin{itemize}[leftmargin=*]
    \item if the link is far away, the constraint is loose;
    \item if the link is near the obstacle, the constraint becomes active;
    \item at the safety boundary, the robot may not move further into the
    obstacle.
\end{itemize}

\section{How OSCBF Avoids Link Collisions From Task-Space Commands}

This is the most important point for our project.

A common but unsafe pipeline is:

\[
\text{plan a safe end-effector path}
\quad \rightarrow \quad
\text{run inverse kinematics}
\quad \rightarrow \quad
\text{execute joint path}.
\]

The failure is that ``safe end-effector path'' only means the tool point is
safe. It says nothing about the rest of the robot. An IK solution may place the
elbow, wrist, or forearm inside an obstacle even when the end effector is far
from it.

OSCBF avoids this by putting the end-effector task and the whole-body safety
constraints into the same optimization problem.

\subsection{Task-Space Motion Is The Objective}

Suppose the task-space planner or density controller gives a desired
end-effector velocity:

\[
\dot{x}_{\mathrm{nom}}.
\]

The QP tries to choose a joint velocity $\dot{q}$ that realizes this motion:

\[
J_v(q)\dot{q} \approx \dot{x}_{\mathrm{nom}}.
\]

So the objective contains:

\[
\|J_v(q)\dot{q} - \dot{x}_{\mathrm{nom}}\|_2^2.
\]

This says:

\begin{quote}
Move the end effector the way the task-space controller wants.
\end{quote}

\subsection{Link Safety Is A Constraint}

For every robot body sphere $i$ and obstacle sphere $j$, define:

\[
h_{ij}(q)
= \|p_i(q) - p_{j,\mathrm{obs}}\|_2
- r_i
- r_{j,\mathrm{obs}}
- d_{\mathrm{safe}}.
\]

This is not just an end-effector distance. The point $p_i(q)$ can be on any
link:

\[
p_i(q) \in
\{\text{base link, shoulder link, elbow link, forearm link, wrist link, tool}\}.
\]

For each such pair, OSCBF adds:

\[
\nabla h_{ij}(q)^\top \dot{q}
\ge
-\gamma h_{ij}(q).
\]

This says:

\begin{quote}
Do not choose a joint velocity that makes any robot link move into an obstacle
too quickly.
\end{quote}

At the boundary of safety, where $h_{ij}(q)=0$, the condition becomes:

\[
\nabla h_{ij}(q)^\top \dot{q} \ge 0.
\]

So the robot cannot move that link any deeper toward the obstacle.

\subsection{The Combined QP}

The combined velocity-control problem is:

\[
\begin{aligned}
\dot{q}^{\star}
= \arg\min_{\dot{q}} \quad
& \|J_v(q)\dot{q} - \dot{x}_{\mathrm{nom}}\|_2^2
 + \lambda \|\dot{q} - \dot{q}_{\mathrm{posture}}\|_2^2 \\
\text{s.t.} \quad
& \nabla h_{ij}(q)^\top \dot{q}
  \ge -\gamma h_{ij}(q),
  \qquad \forall\; \text{robot body } i,\; \text{obstacle } j,\\
& \dot{q}_{\min} \le \dot{q} \le \dot{q}_{\max}.
\end{aligned}
\]

The objective is task-space:

\[
\text{track the desired end-effector motion}.
\]

The constraints are whole-body:

\[
\text{keep every modeled link away from obstacles}.
\]

This is the key separation. Task space gives the desired behavior. The CBF
constraints decide which joint velocities are allowed.

\subsection{Why This Fixes The IK Failure Mode}

In the unsafe IK pipeline, collision checking happens after the task-space
motion has already been converted into a joint path. If the IK solution is bad,
we discover the problem late.

In OSCBF, the optimizer never treats IK as a separate final step. It directly
chooses $\dot{q}$ while enforcing link-safety constraints.

So instead of:

\[
\text{safe tool path} \rightarrow \text{possibly unsafe IK},
\]

we solve:

\[
\text{best joint velocity for the tool task}
\quad
\text{subject to all link distances staying safe}.
\]

If the nominal end-effector command would require the elbow to move into a
cylinder, the link-obstacle CBF for the elbow becomes active and the QP finds
another feasible joint velocity. If no safe joint velocity exists, the QP
slows, stops, or uses slack depending on the implementation.

\section{Joint Limit Barrier}

For each joint:

\[
h_{\min,i}(q) = q_i - q_{\min,i},
\]

\[
h_{\max,i}(q) = q_{\max,i} - q_i.
\]

The CBF constraints are:

\[
\dot{q}_i \ge -\gamma(q_i - q_{\min,i}),
\]

\[
-\dot{q}_i \ge -\gamma(q_{\max,i} - q_i).
\]

These slow the joint as it approaches its limits.

\section{Singularity Avoidance Barrier}

The paper also considers singularity avoidance using manipulability:

\[
h(q) = \mu(q) - \epsilon.
\]

The manipulability $\mu(q)$ can be written using singular values of the
Jacobian:

\[
\mu(q) = \prod_i \sigma_i(J(q)).
\]

Safe means:

\[
\mu(q) \ge \epsilon.
\]

This keeps the robot away from configurations where the Jacobian becomes close
to singular.

\section{Why This Is Nice}

The formulation is attractive for manipulation because:

\begin{itemize}[leftmargin=*]
    \item It handles whole-body safety, not only end-effector safety.
    \item It can enforce many constraints at the same time.
    \item The QP is local and fast compared with nonlinear MPC.
    \item It is less heuristic than artificial potential fields.
    \item The safety filter is minimally invasive.
    \item The objective is task-consistent, so the robot preserves the
    end-effector task instead of merely staying close to an arbitrary joint
    command.
\end{itemize}

\section{Velocity, Torque, And Position Control}

The OSCBF GitHub repository is organized around two main controller
formulations:

\begin{itemize}[leftmargin=*]
    \item \textbf{velocity control}: the safety filter outputs joint velocities
    $\dot{q}$;
    \item \textbf{torque control}: the safety filter outputs joint torques
    $\tau$.
\end{itemize}

The repository README says OSCBF can enforce safety on both torque-controlled
and velocity-controlled robots, and can be used either as a full controller or
as a safety filter on top of an existing controller. The code reflects this:
there are separate OSCBF configuration classes for velocity control and torque
control.

\subsection{Velocity-Controlled OSCBF}

In velocity control, the state and input are:

\[
z = q,
\qquad
u = \dot{q}.
\]

The model is:

\[
\dot{q} = u.
\]

For a safety function $h(q)$, the derivative is:

\[
\dot{h}(q, \dot{q})
= \nabla h(q)^\top \dot{q}.
\]

The CBF constraint is:

\[
\nabla h(q)^\top \dot{q}
\ge
-\gamma h(q).
\]

This is the cleanest version for our first implementation because the QP
directly chooses the safe joint velocity.

\subsection{Torque-Controlled OSCBF}

In torque control, the state and input are:

\[
z =
\begin{bmatrix}
q \\
\dot{q}
\end{bmatrix},
\qquad
u = \tau.
\]

The dynamics are second order:

\[
M(q)\ddot{q} + c(q,\dot{q}) + g(q) = \tau.
\]

Equivalently:

\[
\begin{bmatrix}
\dot{q} \\
\ddot{q}
\end{bmatrix}
=
f(q,\dot{q}) + g(q)\tau.
\]

For collision avoidance, joint limits, and workspace limits, the safety
function often depends on position:

\[
h = h(q).
\]

The torque $\tau$ does not appear in $\dot{h}$ immediately. It appears after
differentiating again:

\[
\ddot{h}(q,\dot{q},\tau)
= L_f^2 h(z) + L_g L_f h(z)\tau.
\]

So torque control usually needs a high-order CBF:

\[
L_f h_2(z) + L_g h_2(z)\tau
\ge
-\alpha_2(h_2(z)),
\]

where:

\[
h_2(z) = L_f h(z) + \alpha_1(h(z)).
\]

This is more physically accurate, especially for fast motions, because it
accounts for inertia, Coriolis/centrifugal effects, gravity, and torque limits.
But it is also more complex.

\subsection{Where Position Control Fits}

Position control is different. A position-controlled robot accepts a desired
joint position:

\[
q_{\mathrm{cmd}}.
\]

It does not directly accept the CBF decision variable $\dot{q}$ or $\tau$.
Instead, an internal controller tries to move the robot toward
$q_{\mathrm{cmd}}$.

For example, the low-level behavior may look roughly like:

\[
\dot{q} \approx K(q_{\mathrm{cmd}} - q).
\]

This means a pure CBF constraint in terms of $\dot{q}$ is not exactly the same
as commanding positions. However, we can still use the velocity OSCBF as an
outer-loop planner:

\[
\dot{q}^{\star}
= \text{OSCBF}(q, \dot{x}_{\mathrm{nom}}, \text{safety constraints}),
\]

then convert it to a small position command:

\[
q_{\mathrm{cmd}}
= q + \Delta t\,\dot{q}^{\star}.
\]

This is the practical bridge to MoveIt/Gazebo trajectory execution. The CBF
thinks in velocities, while the simulator/controller receives short position
waypoints.

\subsection{Recommended Choice For UF850}

For our first UF850 simulation:

\begin{itemize}[leftmargin=*]
    \item use the \textbf{velocity-controlled OSCBF formulation};
    \item compute $\dot{q}^{\star}$ from the QP;
    \item integrate it into small position waypoints;
    \item send those waypoints through the existing position/trajectory
    controller.
\end{itemize}

This gives us the simple CBF math:

\[
\nabla h(q)^\top \dot{q}
\ge
-\gamma h(q),
\]

while remaining compatible with MoveIt-style execution:

\[
q_{\mathrm{next}} = q + \Delta t\,\dot{q}^{\star}.
\]

Later, if we need faster and more dynamically faithful behavior, we can move to
the torque-controlled OSCBF formulation.

\section{How We Should Use It For UF850}

For our first density-based UF850 planner:

\[
\dot{x}_{\mathrm{nom}}
= u_{\mathrm{density}}(x_{\mathrm{tcp}}, x_{\mathrm{goal}}, \text{workspace obstacles}).
\]

Then solve:

\[
\begin{aligned}
\dot{q}^{\star}
= \arg\min_{\dot{q}} \quad
& \|J_v(q)\dot{q} - \dot{x}_{\mathrm{nom}}\|_2^2
 + \lambda \|\dot{q} - \dot{q}_{\mathrm{posture}}\|_2^2 \\
\text{s.t.} \quad
& \nabla h_{\mathrm{collision},k}(q)^\top \dot{q}
  \ge -\gamma h_{\mathrm{collision},k}(q), \\
& \nabla h_{\mathrm{joint},k}(q)^\top \dot{q}
  \ge -\gamma h_{\mathrm{joint},k}(q), \\
& \dot{q}_{\min} \le \dot{q} \le \dot{q}_{\max}.
\end{aligned}
\]

Then integrate:

\[
q_{t+1} = q_t + \Delta t \, \dot{q}^{\star}.
\]

This replaces the risky pipeline:

\[
\text{safe end-effector path} \rightarrow \text{IK} \rightarrow
\text{possible link collision}
\]

with:

\[
\text{desired end-effector velocity} \rightarrow
\text{whole-arm safe QP} \rightarrow
\text{safe joint velocity}.
\]

\section{Takeaway}

The paper's key message is:

\begin{quote}
When a robot's task is defined in operational space, its safety filter should
also preserve operational-space task behavior while enforcing whole-body safety.
\end{quote}

For our project, that means density feedback should generate the desired tool
motion, and OSCBF-style constraints should make that motion safe for all UF850
links and joints.

\end{document}
